\documentclass{article}

\usepackage[final]{neurips_2024}
\usepackage[utf8]{inputenc}
\usepackage[T1]{fontenc}
\usepackage{amsmath}
\usepackage{amssymb}
\usepackage{booktabs}
\usepackage{graphicx}
\usepackage{subcaption}
\usepackage{microtype}
\usepackage{hyperref}

\setlength{\parskip}{3pt}
\graphicspath{{../figures/signal_analysis/}{figures/signal_analysis/}}

\title{Measuring Hawkish and Dovish Signals in FOMC Communications\\
{\large A Text-as-Data Analysis of Statements and Minutes, 2000--2024}}

\author{Mathieu Janin\\
ENSAE Paris\\
Machine Learning for NLP}

\begin{document}

\maketitle

\begin{abstract}

Central bank communication is a key channel through which monetary policy affects expectations beyond actual interest-rate decisions. This paper studies how the Federal Open Market Committee (FOMC) 
conveys hawkish and dovish signals in its \textit{Statements} and \textit{Minutes} between 2000 and 2024. We construct a domain-specific dictionary that contrasts terms associated with inflation pressures, 
tightening, and upside risks with terms related to easing, downside risks, and labor-market weakness. The resulting tone score is computed at both document and sentence level, and is evaluated through 
several robustness checks, including the masking of explicit rate-decision sentences, dictionary sensitivity tests, TF-IDF classification benchmarks, latent semantic analysis, and Sentence-BERT embeddings.
The results show that hawkish language is substantially stronger during rate-hike meetings, while high-volatility periods are associated with a more dovish tone. By contrast, meetings around U.S. elections 
do not display systematic differences. \textit{Statements} and \textit{Minutes} have similar average tone, but some meetings exhibit large document-level gaps, suggesting that \textit{Minutes} often add 
nuance to the immediate policy message. Finally, embedding-based representations capture broader semantic and macroeconomic structure, but they do not reproduce the dictionary-based hawkish/dovish axis 
without supervision. Overall, the results support the use of interpretable, domain-specific lexical measures for tracking monetary-policy tone, while highlighting the limits of generic embeddings for 
measuring such a targeted economic concept.

\end{abstract}

\section{Introduction}

The mandate of the Federal Reserve is dual: it must contain inflation while maintaining a dynamic labor market. To achieve this, it has at its disposal monetary policy tools, notably the management of policy interest rates. 
Its mission also consists of describing economic conditions, signaling risks, and guiding expectations regarding the future path of monetary policy. This communication channel is particularly visible at the conclusion of meetings of 
the Federal Open Market Committee (FOMC), when the Committee releases a short Statement, followed several weeks later by more detailed records (\textit{Minutes}). These documents are official, published on a 
recurring basis, and closely analyzed by financial markets. They therefore constitute a natural corpus for studying monetary policy communication using NLP methods.
This article poses a \textit{text-as-data} question: to what extent do FOMC communications convey a ``hawkish'' or ``dovish'' signal, and how does this signal vary depending on the type of 
document (\textit{Statements} / \textit{Minutes}), rate decisions (\textit{Cut} / \textit{Hold} / \textit{Hike}), financial market volatility, electoral periods, and macroeconomic regimes? A ``hawkish'' signal emphasizes inflation, 
price pressures, restrictive financial conditions, monetary tightening, or higher interest rates. Conversely, a ``dovish'' signal highlights support for economic activity, labor market weakness, monetary easing, lower interest rates, 
or downside macroeconomic risks.

The object of interest (the ``hawkish''/``dovish'' signal) is latent: it is not directly observed in the data and must be inferred from words, expressions, and the structure of documents. We use a dictionary-based measure as the 
main estimator and compare it with TF-IDF, LSA, and Sentence-BERT representations. The dictionary score constitutes the central measure, as it corresponds transparently to well-defined economic concepts. 
TF-IDF classifiers assess whether information related to rate decisions can also be extracted from a more flexible \textit{bag-of-words} representation. Semantic representations derived from latent semantic analysis (LSA) 
and Sentence-BERT, in turn, examine whether a dense semantic geometry produces a similar ``hawkish/dovish'' ranking or instead captures a broader institutional and macroeconomic structure.
The results presented in this article are descriptive rather than causal. They show that meetings associated with policy rate increases are clearly more ``hawkish'' than those with no change or with rate cuts (as expected), 
and that meetings characterized by a high level of financial volatility, measured by the VIX index, are associated with a less ``hawkish'' tone (also expected, since volatility is negatively correlated with market returns; thus, 
periods of crisis or weak growth tend to be associated with higher volatility). Electoral periods do not exhibit a marked average difference. The \textit{Statements} and the \textit{Minutes} display a similar average tone, 
but differences at the level of individual meetings can be substantial, suggesting that the minutes often contextualize or nuance the immediate message released after the meeting. Finally, combining the ``hawkish/dovish'' score 
with an inflation--employment balance indicator shows that the economic diagnosis and the monetary policy signal are closely linked—particularly in the \textit{Minutes}—without constituting the same textual object.
The paper makes three contributions. First, it constructs a reproducible processing pipeline that systematically aligns FOMC \textit{Statements} and \textit{Minutes} by meeting and enriches them with financial, 
political, and macroeconomic context. Second, it proposes a transparent and interpretable measure of monetary policy tone that can be applied at both the document and sentence levels. Third, it compares this manually constructed 
measure with more flexible text representations derived from machine learning. This comparison separates two tasks that are often conflated: measuring a predefined economic concept and predicting rate-decision labels from text.

\section{Related Work}
Over the past decades, central bank communication has become an instrument of monetary policy in its own right. \citet{blinder2008communication} review the theoretical and empirical evidence and argue that communication 
can make policy more predictable, influence asset prices, and shape expectations. This literature has helped establish the idea that FOMC texts can be analyzed as economic data rather than as a simple institutional narrative.
Several articles use computational methods to study FOMC communication. \citet{lucca2009measuring} construct automated measures of the content of FOMC statements and show that long-term interest rates respond to communication surprises. 
This result is particularly relevant for our study, as the same rate decision can be accompanied by different language regarding the future orientation of monetary policy. \citet{hansen2018transparency} apply NLP methods to FOMC 
transcripts and show that transparency alters patterns of deliberation and communication. Our article instead focuses on Statements and Minutes, which are regularly available and comparable over time, but it is grounded in the same 
perspective that monetary communication can be measured systematically.
The article also fits within the broader \textit{text-as-data} literature. \citet{gentzkow2019text} emphasize that textual analysis requires explicit choices regarding representation, measurement, 
and validation. \citet{loughran2011liability} show that generic sentiment dictionaries can be misleading in financial corpora, where common words may have domain-specific meanings. For this reason, the main measure proposed here 
is not an index of positive/negative sentiment, but rather an index of monetary policy tone. The embedding-based checks rely on latent semantic analysis \citep{deerwester1990indexing} and on 
Sentence-BERT \citep{reimers2019sentencebert}, which provide alternative semantic representations of the same documents.
This article is methodologically close to dictionary-based approaches, while using machine learning representations as diagnostic tools rather than as the primary estimator. This choice reflects an empirical constraint specific 
to the corpus. The sample is relatively small by current NLP standards: after alignment, the unit of analysis comprises fewer than two hundred meetings. In the absence of expert-labeled sentences, a fully supervised ``hawkish/dovish'' 
classifier would risk learning language specific to particular periods or document formats rather than the underlying monetary policy signal. A dictionary is admittedly less flexible, but it provides a direct and interpretable link 
between the measured score and the economic vocabulary used to construct it.

\section{Data}

\subsection{FOMC Corpus}

The corpus contains FOMC statements and minutes collected from the Federal Reserve
website \citep{federalreservefomc}. Statements are the short announcements
released immediately after each meeting. Minutes are longer documents published
later; they summarize the discussion, the uncertainty around the outlook, and the
arguments behind the decision.

The final aligned sample covers 199 meetings from February 2000 to December 2024.
Each observation contains a statement, the corresponding minutes, the meeting date
and contextual variables. Rate decisions are inferred from the statement text and
classified as \emph{hike}, \emph{hold} or \emph{cut}. The sample contains 59 hikes,
110 holds and 30 cuts.

\begin{table}[t]
\centering
\caption{Aligned corpus description.}
\label{tab:corpus}
\begin{tabular}{lrrrrr}
\toprule
Document type & Meetings & Mean words & Median words & Mean sentences & Date range \\
\midrule
Statements & 199 & 526.1 & 536.0 & 26.8 & 2000--2024 \\
Minutes & 199 & 7282.7 & 7518.0 & 311.9 & 2000--2024 \\
\bottomrule
\end{tabular}
\end{table}

Table \ref{tab:corpus} shows the large difference in document length. Minutes are
roughly fourteen times longer than statements on average. This makes length
normalization essential and also explains why statements and minutes should not be
expected to carry the same kind of signal. Statements are calibrated policy
messages while minutes are richer institutional summaries.

\subsection{Alignment and Preprocessing}

The raw corpus is organized around the dates of FOMC meetings. Statements and meeting records, i.e. Minutes, are stored in separate fields and then systematically aligned by meeting date. 
Preprocessing is deliberately conservative in order to preserve as much of the original textual information as possible. Text is converted to lowercase for counting purposes, whitespace is normalized, and recurring fragments 
originating from web pages (such as navigation menus, update notices, and accessibility statements) are removed when they appear. These steps remove recurring webpage noise while preserving monetary-policy content.
The analysis does not remove stopwords that are economically meaningful in the computation of dictionary scores, because several compound expressions depend on local syntax, for example \textit{higher rates}, \textit{downside risks}, 
and \textit{price stability}. Stopword removal is used only in the TF-IDF benchmark, where the objective is predictive comparison rather than direct term counting. This distinction reflects the conceptual difference between an 
interpretable dictionary-based measure and a statistical representation designed to maximize classification performance.
Documents are also segmented into sentences. Sentence segmentation serves two main functions. First, it allows the masking of sentences that explicitly announce the rate decision, in order to prevent the textual score from 
mechanically reflecting monetary policy information that is already observed. Second, it enables a dictionary audit at the sentence level by precisely identifying the text segments that contribute to the final score. 
Sentence-level labels guard against scores driven by a few repeated terms and make the document-level measure easier to audit.

\begin{table}[t]
\centering
\caption{Rate-decision distribution in the aligned sample.}
\label{tab:decisions_sample}
\begin{tabular}{lrr}
\toprule
Decision class & Meetings & Share \\
\midrule
Cut & 30 & 15.1\% \\
Hold & 110 & 55.3\% \\
Hike & 59 & 29.6\% \\
\bottomrule
\end{tabular}
\end{table}

Table \ref{tab:decisions_sample} reports the class balance. The dominance of
hold meetings is expected in monetary policy data, but it matters for evaluation:
accuracy alone can be misleading. This is why the predictive experiments also
report macro-F1, which gives equal weight to the three decision classes.

\subsection{Contextual Variables}

In our empirical analysis, we introduce three sets of contextual variables. First, daily data on the VIX volatility index are obtained from Cboe. For each FOMC meeting, the VIX is averaged over a window ranging from 
five days before to five days after the meeting; the upper quartile of this distribution defines meetings characterized by high financial volatility. This measure is intended to capture episodes of market stress that may influence 
the tone of monetary communication independently of rate decisions.
Second, electoral periods are defined as meetings held within the 90 days preceding a U.S. presidential election or midterm elections. This 
variable tests whether policy language changes near politically sensitive dates, without assuming strategic behavior.

Third, meetings are also classified according to broad macro-financial regimes corresponding to distinct phases of monetary policy and the macroeconomic environment. These regimes include: (i) the early-2000s easing cycle, marked 
by the slowdown following the burst of the dot-com bubble; (ii) the global financial crisis of 2008--2009, characterized by systemic financial instability and emergency policy measures; (iii) the zero lower bound (\textit{ZLB}) period, 
during which policy rates remained close to zero and monetary policy relied more heavily on unconventional instruments; (iv) the normalization phase of 2015--2019, associated with a gradual recovery in economic activity and a 
progressive tightening of monetary policy; (v) the Covid-19 crisis in 2020, marked by a sharp economic shock and exceptional monetary interventions; (vi) the post-Covid inflation period beginning in 2021, characterized by persistent 
inflation acceleration and a rapid cycle of monetary tightening; and (vii) a reference category grouping other meetings that do not fall explicitly within these major episodes.
These contextual variables are not used to identify causal effects. Rather, they serve to structure the descriptive analysis and to verify whether the textual signal evolves in an economically coherent manner. For example, a 
meaningful measure of tone should be more ``hawkish'' during rate hikes than during rate cuts, but it should also reveal more nuanced patterns during periods of stress, when the Federal Reserve may refer to inflationary risks while 
simultaneously emphasizing financial fragility or downside macroeconomic risks.

\section{Methodology and Results}

\subsection{Dictionary Tone Score}

The main model is a dictionary score. Let $H_{d,t}$ and $D_{d,t}$
denote the length-normalized counts of hawkish and dovish terms in document
$d$ for meeting $t$:
\begin{equation}
H_{d,t} = \frac{\#\text{hawkish terms}_{d,t}}{\#\text{words}_{d,t}},
\qquad
D_{d,t} = \frac{\#\text{dovish terms}_{d,t}}{\#\text{words}_{d,t}}.
\end{equation}
The normalized tone score is
\begin{equation}
\text{ToneNorm}_{d,t} =
\frac{H_{d,t}-D_{d,t}}{H_{d,t}+D_{d,t}}.
\end{equation}
Positive values indicate a more hawkish document, negative values indicate a more
dovish document, and values close to zero indicate a balanced document.

The hawkish dictionary includes terms such as \emph{inflation},
\emph{price stability}, \emph{elevated}, \emph{tightening},
\emph{restrictive}, \emph{higher rates}, \emph{persistent inflation}, and
\emph{inflationary pressures}. The dovish dictionary includes terms such as
\emph{unemployment}, \emph{weakness}, \emph{slowdown}, \emph{accommodative},
\emph{support}, \emph{lower rates}, \emph{easing}, \emph{recovery}, and
\emph{downside risks}. The score is intentionally interpretable: its changes can
be audited by reading the terms and the sentences that drive them.

The normalization in Equation (2) is chosen because it is scale-free. A raw
difference $H_{d,t}-D_{d,t}$ would still depend on the absolute density of
dictionary terms. The normalized balance instead asks whether policy-related
language is more hawkish or more dovish conditional on the document containing
such language. This is particularly useful when comparing statements and minutes,
because minutes contain far more words and far more policy-relevant sentences.
When both $H_{d,t}$ and $D_{d,t}$ are zero, the score is set to zero; in practice
this is rare for FOMC texts.

\subsection{Choice of Model Family}

The dictionary score is selected as the primary model because the research question is explicitly oriented toward measurement rather than prediction. The article aims to construct an interpretable time series of monetary policy 
tone that is comparable across meetings and across document types. For this objective, a transparent and stable measure is preferable to a ``black-box'' classifier trained on a relatively limited sample.
The additional natural language processing models address distinct methodological questions. TF-IDF, combined with linear models, tests whether the text contains sufficient information to predict the observed rate decision, 
independently of the interpretable structure of the dictionary. Latent semantic analysis (LSA) examines whether documents that are semantically close in a reduced \textit{bag-of-words} space also appear similar along the dimension 
of monetary policy tone. Sentence-BERT evaluates whether pre-trained contextual embeddings, without task-specific adjustment (\textit{fine-tuning}), spontaneously align with the same ``hawkish/dovish'' prototype contrast.
These approaches are therefore complementary rather than competing. The dictionary is designed to measure an economic concept defined \textit{a priori}, while TF-IDF and embeddings serve as validation and diagnostic tools, making 
it possible to assess the robustness and coherence of this measure under alternative representations of the corpus.

\subsection{Validation and Robustness Checks}

Four main checks are implemented to assess whether the dictionary captures an economically meaningful monetary policy signal.
First, a masked score removes sentences that explicitly announce the rate decision. This procedure makes it possible to test whether the measured tone merely restates the decision—for example, 
``the Committee raised the target range''—or whether it also appears in the economic assessment and justification surrounding that decision. The objective is to prevent the measure from mechanically reflecting information that 
is already observed.
Second, a sentence-level version assigns each sentence a label—``hawkish,'' ``dovish,'' mixed, or neutral—based on dictionary occurrences. This approach provides both a qualitative analysis of document content and a 
quantitative measure, at the document level, of the proportion of sentences relevant to monetary policy. It also strengthens the interpretability of the score by explicitly identifying the textual segments that contribute to 
the measured tone.
Third, TF-IDF benchmarks evaluate whether the rate decision can be predicted from the text. Logistic regression and linear SVM models are trained using chronological splits, with earlier meetings used for training and later 
meetings used for out-of-sample evaluation. This temporal split is more demanding than a random split, because the language of monetary policy evolves over time. Models are estimated on full texts, texts masking the decision, 
and texts stripped of standardized elements (\textit{boilerplate}), for both \textit{Statements} and \textit{Minutes}. The TF-IDF representation uses unigrams and bigrams, stopword removal, a minimum document frequency of two, and 
a maximum vocabulary size of 6\,000 features. The test set consists of the most recent 30\% of meetings, covering the period from June 2017 to December 2024, while the training set covers the period from February 2000 to May 2017.
Fourth, dense representations are used as semantic robustness checks. Latent semantic analysis (LSA) is computed by applying truncated singular value decomposition (SVD) to a TF-IDF matrix of document excerpts. A tone score is 
then defined as the difference between cosine similarity with a ``hawkish'' prototype and cosine similarity with a ``dovish'' prototype. Sentence-BERT applies the same prototype logic to contextual embeddings generated by 
the \texttt{all-MiniLM-L6-v2} model. These embedding-based scores serve as robustness checks to determine whether a semantic space less constrained by manual construction produces a comparable ranking of documents.
Finally, the study uses a pooled ordinary least squares (\textit{pooled OLS}) regression model as a synthetic descriptive check. The dependent variable is the tone score for each document. Explanatory variables include 
document type, indicator variables for rate decisions, the standardized VIX, and an electoral period indicator. This model is a descriptive tool intended to verify whether the associations observed in the figures remain present 
when these different dimensions are taken into account simultaneously.

\begin{figure}[t]
\centering
\includegraphics[width=0.93\textwidth]{tone_norm_over_time.png}
\caption{Normalized hawkish/dovish tone in FOMC statements and minutes.}
\label{fig:tone_time}
\end{figure}

Figure \ref{fig:tone_time} shows that the average tone is positive in both
document types. The mean normalized score is 0.484 for statements and 0.507 for
minutes. This means that, over the full sample, language related to inflation, price pressure and restrictive
conditions is more frequent than language related to easing or weakness. The score
falls in the early 2000s and in crisis periods, then rises during the 2015--2019
normalization and the post-Covid inflation episode.

The masked score is very close to the main score: the means are 0.474 for
statements and 0.501 for minutes. The policy signal is therefore not driven only
by direct decision sentences. The same gradient appears at sentence level:
policy-relevant sentences account for about 31.0\% of statements sentences and
29.7\% of minutes sentences, and their balance is highest during hike meetings.

The time series shows why average tone alone is insufficient. During
the zero lower bound period, policy was highly accommodative, yet the documents
continued to discuss price stability and inflation expectations. During the
post-Covid inflation period, the language becomes more directly hawkish because
the inflation diagnosis and the tightening signal point in the same direction.
The dictionary therefore behaves less like a generic sentiment index and more like
a stance indicator: it reacts to the mixture of inflation-risk language and
accommodation language, not to whether the economic outlook is described as
``good'' or ``bad''.

\subsection{Statements, Minutes and Rate Decisions}

\begin{figure}[t]
\centering
\begin{subfigure}{0.48\textwidth}
\centering
\includegraphics[width=\textwidth]{tone_by_rate_decision.png}
\caption{Tone by decision.}
\end{subfigure}
\hfill
\begin{subfigure}{0.48\textwidth}
\centering
\includegraphics[width=\textwidth]{statement_minutes_tone_distance.png}
\caption{Statements-minutes distance.}
\end{subfigure}
\caption{Document tone by rate decision and meeting-level distance between
statements and minutes.}
\label{fig:decision_distance}
\end{figure}

Hike meetings are clearly more hawkish than hold and cut meetings. Table
\ref{tab:decision} reports average scores. For statements, hike meetings have an
average tone of 0.705, compared with 0.384 for holds and 0.418 for cuts. For
minutes, the corresponding values are 0.700, 0.442 and 0.368. The lower average
score for cut minutes is economically intuitive: minutes written around easing
episodes emphasize weakness, downside risks and support more than hike minutes do.

\begin{table}[t]
\centering
\caption{Average normalized tone by rate decision.}
\label{tab:decision}
\begin{tabular}{lrrr}
\toprule
Document type & Cut & Hold & Hike \\
\midrule
Statements & 0.418 & 0.384 & 0.705 \\
Minutes & 0.368 & 0.442 & 0.700 \\
\bottomrule
\end{tabular}
\end{table}

The statements-minutes distance is
\begin{equation}
\text{DistanceTone}_{t} =
|\text{ToneNorm}_{minutes,t}-\text{ToneNorm}_{statements,t}|.
\end{equation}
Its mean is 0.163 and its median is 0.122. A paired comparison does not indicate a
large systematic difference in average tone between statements and minutes, but
the meeting-level distribution is informative. Several meetings, especially in
the early-2000s easing cycle, display large gaps. In those cases, the Statement
can be strongly dovish while the minutes remain more balanced because they record
competing risks or a more detailed debate.

Statements and Minutes differ not only in length but also in function. Statements are written as coordinated
public signals. They must summarize the Committee's decision in a compact and
market-readable form. Minutes, by contrast, are designed to document the reasoning
behind the decision. They can contain hawkish and dovish arguments in the same
document, especially when the Committee faces a trade-off between inflation
concerns and real-activity weakness. A large statements-minutes distance should
therefore not be interpreted automatically as inconsistency. It is better read as
evidence that the later document adds nuance to the immediate signal.

\subsection{Financial and Political Context}

\begin{figure}[t]
\centering
\begin{subfigure}{0.48\textwidth}
\centering
\includegraphics[width=\textwidth]{tone_and_vix_over_time.png}
\caption{Tone and VIX over time.}
\end{subfigure}
\hfill
\begin{subfigure}{0.48\textwidth}
\centering
\includegraphics[width=\textwidth]{tone_by_macro_regime.png}
\caption{Tone by macro regime.}
\end{subfigure}
\caption{Tone across financial stress and macro-financial regimes.}
\label{fig:context}
\end{figure}

High-volatility meetings are associated with lower tone scores. In meetings below
the upper VIX quartile, the average score is 0.542 for statements and 0.549 for
minutes. In high-VIX meetings, the averages fall to 0.311 and 0.382. This pattern
is consistent with the Fed using more dovish or support-oriented language during
episodes of financial stress. In contrast, election windows show little average
difference: statement tone is 0.494 outside election windows and 0.459 inside
them; tone in Minutes is 0.511 outside and 0.497 inside.

The macro-regime analysis is consistent with the time series. The normalization
period and the post-Covid inflation period are among the most hawkish regimes.
The Covid period and the early-2000s easing cycle are less hawkish. The global
financial crisis is not the most dovish regime in the dictionary score because
those documents also contain frequent references to inflation, price stability and
financial strains. The global financial crisis shows that weakness language and 
inflation-risk language can coexist in the same document.

The absence of a strong election-window pattern is also informative. In this
descriptive design, the tone series does not suggest an obvious shift in the
language of policy around presidential or midterm elections. This should not be
overinterpreted as evidence that politics never affects communication. The
election variable is deliberately simple, and the number of meetings inside
election windows is limited. The result is narrower: conditional on this corpus
and this dictionary score, the most visible variation is linked to policy and
financial conditions, not to the election-window indicator.

\subsection{Regression Evidence}

\begin{table}[t]
\centering
\caption{Descriptive OLS model for normalized tone. The reference category is a
statement during a hold meeting.}
\label{tab:ols}
\begin{tabular}{lrr}
\toprule
Variable & Coefficient & $p$-value \\
\midrule
Intercept & 0.406 & 0.000 \\
Minutes & 0.023 & 0.376 \\
Hike decision & 0.255 & 0.000 \\
Cut decision & 0.012 & 0.755 \\
VIX, standardized & -0.082 & 0.000 \\
Election window & 0.002 & 0.945 \\
\bottomrule
\end{tabular}
\end{table}

Table \ref{tab:ols} summarizes the descriptive associations in a pooled document
panel. The hike coefficient is large and statistically significant, which confirms
that the tone score is aligned with the most tightening-oriented decision class.
The VIX coefficient is negative and significant. The minutes dummy is positive but
not statistically significant, consistent with the paired comparison. The election
coefficient is essentially zero in this specification.

The cut coefficient is small in the pooled specification, even though cut minutes
are visibly less hawkish in the grouped averages. The small cut coefficient reflects 
heterogeneity across easing episodes. Some cuts occur in crisis periods dominated by
support and downside-risk language; others occur when inflation concerns have not
fully disappeared. The tone score therefore captures the communication
surrounding the decision, not only the sign of the rate change.

\subsection{NLP Model Checks}

\begin{figure}[t]
\centering
\begin{subfigure}{0.48\textwidth}
\centering
\includegraphics[width=\textwidth]{tfidf_policy_models_macro_f1.png}
\caption{TF-IDF chronological test macro-F1.}
\end{subfigure}
\hfill
\begin{subfigure}{0.48\textwidth}
\centering
\includegraphics[width=\textwidth]{dictionary_sensitivity.png}
\caption{Main and strict dictionary scores.}
\end{subfigure}
\caption{Representation and dictionary robustness checks.}
\label{fig:model_checks}
\end{figure}

The TF-IDF models provide a benchmark for whether rate-decision information is
present in the text beyond the handcrafted score. The best chronological test
performance is obtained with minutes and a linear SVM: 0.550 accuracy and 0.407
macro-F1. The same model on masked minutes reaches the same accuracy and 0.401
macro-F1, which again suggests that the decision signal is not confined to the
explicit announcement sentence. Models trained on Statements perform worse, with the best
macro-F1 around 0.333. This gap is plausible because statements are short and
formulaic, while minutes contain more evidence about the economic reasoning behind
the decision.

The TF-IDF coefficients are useful for auditing the model. Hike-associated terms
include inflation, prices, energy prices and point increases. Cut-associated terms
include discount-rate reductions and crisis-era action vocabulary. Hold-associated
terms include purchases, recovery, accommodative stance and asset programs. The
terms are economically interpretable, but they also reveal a limitation: a flexible
text classifier can partly learn episode-specific vocabulary, not only policy
stance.

Dictionary sensitivity is checked by removing broad words such as generic
\emph{increase}, \emph{pressure} and \emph{support}. The main and strict scores
remain strongly related. The correlation is 0.839 for statements and 0.935 for
minutes. The mean absolute difference is larger for statements (0.159) than for
minutes (0.079), which is expected because short documents are more sensitive to
small dictionary changes. The ordering by rate decision remains stable.

\begin{table}[t]
\centering
\caption{Summary of representation checks.}
\label{tab:robustness}
\begin{tabular}{p{0.30\textwidth}p{0.24\textwidth}p{0.24\textwidth}}
\toprule
Check & Statements result & Minutes result \\
\midrule
Masked mean tone & 0.474 & 0.501 \\
Strict dictionary correlation with main score & 0.839 & 0.935 \\
Best TF-IDF chronological macro-F1 & 0.333 & 0.407 \\
LSA prototype correlation with dictionary tone & 0.407 & 0.150 \\
Sentence-BERT prototype correlation with dictionary tone & 0.043 & 0.048 \\
\bottomrule
\end{tabular}
\end{table}

Table \ref{tab:robustness} summarizes the main robustness checks. The strict
dictionary result supports the stability of the lexical score. The TF-IDF result
shows that minutes contain more predictive information about the policy decision
than statements do. The embedding correlations are weaker, especially for minutes,
which suggests that dense semantic structure is not equivalent to the explicit
hawkish/dovish concept measured by the dictionary.

\subsection{Embedding Robustness}

\begin{figure}[t]
\centering
\includegraphics[width=0.92\textwidth]{lsa_embedding_robustness.png}
\caption{LSA document geometry and prototype-based tone.}
\label{fig:lsa}
\end{figure}

Figure \ref{fig:lsa} compares the dictionary score with an LSA prototype score.
The first 30 LSA components explain 51.5\% of TF-IDF variance. The correlation
between LSA prototype tone and dictionary tone is 0.407 for statements and 0.150
for minutes. This positive but moderate relationship suggests that LSA captures
some policy-tone information, especially in short statements, while minutes'
dense geometry is shaped by broader topics and document structure.

Nearest-neighbor results lead to a similar interpretation. For selected
meetings, neighboring documents often come from similar policy episodes, such as
crisis meetings, zero lower bound communications, or inflation-tightening periods.
This is useful for qualitative exploration: embeddings can retrieve comparable
communication episodes without relying on the dictionary. However, the neighbors
are not always ordered by hawkishness. They may instead be close because they
share institutional vocabulary, document format or crisis-specific topics. For 
this reason, embeddings are kept as robustness checks rather than used as the main measure.

The Sentence-BERT prototype check produces very weak correlations with the
dictionary tone: 0.043 for statements and 0.048 for minutes. This means that a generic contextual embedding,
combined with short hawkish and dovish prototypes, does not reproduce the explicit
monetary-policy lexicon. In this application, Sentence-BERT is better viewed as a
tool for qualitative nearest-neighbor inspection or clustering of communication
episodes than as a direct substitute for the interpretable dictionary score.

Sentence-BERT embeddings are trained to encode broad sentence similarity, not to estimate central-bank stance. Without
fine-tuning or labelled examples, the model has no reason to place \emph{price
stability}, \emph{restrictive stance}, and \emph{downside risks} on the exact
monetary-policy axis needed here. The weak prototype correlations therefore
support the choice of a domain-specific measure. A supervised embedding model
could be a natural extension, but it would require expert labels at the sentence
or paragraph level.

\subsection{From Economic Diagnosis to Policy Signal}

\begin{figure}[t]
\centering
\includegraphics[width=0.82\textwidth]{mandate_balance_vs_tone.png}
\caption{Inflation/employment balance and hawkish/dovish tone.}
\label{fig:mandate_tone}
\end{figure}

The final analysis relates the hawkish/dovish score to a separate
inflation/employment balance. The correlation between mandate balance and tone is
small for statements (-0.042) but positive for minutes (0.405). The relationship 
between diagnosis and tone is therefore much stronger in Minutes than in Statements. 
Statements are concise and strongly calibrated around the immediate
policy message, so the diagnostic balance and the policy signal can diverge.
Minutes provide more space for connecting inflation, labor-market conditions and
the policy stance; the two dimensions therefore move together more clearly.

Most meetings fall in the inflation-hawkish quadrant. This reflects both the
importance of price stability in FOMC communication and the design of the
hawkish/dovish dictionary. More interesting are the off-diagonal cases: meetings
where inflation language is salient but the tone is relatively dovish, or where
employment concerns coexist with hawkish policy language. These cases are useful
for qualitative interpretation because they identify moments when the diagnosis
and the policy signal do not point in the same direction.

The two-dimensional view separates the Fed’s diagnosis from its policy stance. A document can be
inflation-focused because the Committee is explaining high realized inflation, but
the policy signal can still be cautious if financial conditions are fragile. The
opposite case can occur when labor-market language remains prominent while the
Committee signals that policy must stay restrictive. The scatter plot therefore
turns a one-dimensional tone series into a map of communication regimes.

\section{Discussion of results}

The dictionary score provides a transparent and economically interpretable measure of monetary-policy tone. It reproduces the expected ordering of rate decisions, remains robust to the masking 
of sentences announcing the decision, and remains stable under a stricter version of the dictionary. It also makes it possible to compute interpretable distances, at the meeting level, between \textit{Statements} and \textit{Minutes}.
The results also shed light on the role of document type. The \textit{Statement} is the high-visibility message released immediately after a meeting. It condenses the Committee’s decision and assessment into a short, carefully 
negotiated text. The \textit{Minutes}, by contrast, constitute a delayed explanation of the discussion. The fact that their average tone is not statistically different from that of the statements should not obscure their analytical 
value. Their length makes them more useful for identifying the underlying reasons behind the signal, and their meeting-level distance from the statements helps identify situations in which the public message was subsequently 
qualified or refined through more detailed arguments.
Several limitations remain. Dictionary-based scores depend on the selected terms, and some words are sensitive to context. A term such as ``support'' may describe an accommodative policy, but it may also appear 
in a neutral institutional expression. Conversely, a sentence may be ``hawkish'' without explicitly using any dictionary term. Checks based on a stricter version of the dictionary and on sentence-level analysis reduce this concern, 
but they do not eliminate it entirely.
Another limitation concerns the unit of observation. The analysis treats each document as a single observation, whereas meeting records contain multiple speakers, themes, and sections. A more granular model could distinguish staff 
analyses, participants’ views, and the monetary policy discussion itself. It could also differentiate uncertainty language from policy-stance language. This distinction would be particularly relevant during 
meetings in which the Federal Reserve simultaneously describes weak economic activity, high uncertainty, and inflationary risks.
Supervised checks and embedding-based analyses clarify what the dictionary measures—and what it does not measure. TF-IDF classifiers confirm that information related to monetary policy decisions is indeed present in the text, 
particularly in the meeting records, but they also learn vocabulary specific to certain historical episodes. LSA and Sentence-BERT capture semantic neighborhoods and macroeconomic episodes, but they do not spontaneously reproduce 
the ``hawkish/dovish'' scale without stronger supervision. As an extension, we could manually annotate sentences as ``hawkish,'' ``dovish,'' or neutral, and then train a supervised classifier on expert labels. Another 
extension would be to analyze financial market reactions around meetings in order to distinguish descriptive tone from genuine communication surprises.
The approach may transfer to other central banks, but only after adapting the dictionary to each institution. The same processing pipeline could be applied to European Central Bank press conferences or Bank of England meeting 
records, but the dictionaries would need to be adapted to the institutional language and monetary policy frameworks specific to each institution. Even within the Federal Reserve corpus, vocabulary evolves over time. Terms 
related to asset purchases, forward guidance, and balance sheet management become central after 2008. A static dictionary remains easy to interpret, but future work could allow the vocabulary to evolve in a controlled manner 
while preserving a verifiable correspondence with the ``hawkish'' and ``dovish'' concepts.

\section{Conclusion}

This paper measures the hawkish/dovish signal in FOMC statements and minutes from
2000 to 2024. Hawkish language peaks around rate hikes, declines in high-VIX periods, 
and shows little variation around election windows. Minutes and statements have similar average tone but can differ
substantially at the meeting level. Robustness checks based on decision masking,
sentence tagging, TF-IDF models, strict dictionaries, LSA and Sentence-BERT support
the central interpretation: the dictionary score is a useful, transparent measure
of monetary-policy signal, while dense embeddings are better suited to studying
broader semantic structure and communication regimes.


\begingroup
\scriptsize
\setlength{\bibsep}{0pt}
\bibliographystyle{abbrvnat}
\bibliography{references}
\endgroup

\end{document}